\section*{Hoofdstuk \ref{ch:g12:diversification-of-life} --- De opsplitsing van het leven}

\begin{solution}{exo:g12:diversification-of-life:1}
Genverdubbeling en uiteengroeien; \omterm{prop:g11:antibiotic-resistance:variation}{horizontale genoverdracht}; kruising met
polyploïdie; veranderingen in de regeling van \omterm{prop:g12:diversification-of-life:development}{ontwikkelingsgenen};
\omterm{def:g12:diversification-of-life:symbiosis}{symbiose}; doorgegeven gedrag.
\end{solution}

\begin{solution}{exo:g12:diversification-of-life:2}
Een groep verwante \omterm{def:g10:universal-dna:gene}{genen} die door opeenvolgende verdubbelingen van één
voorouderlijk \omterm{def:g10:universal-dna:gene}{gen} afstammen, waarbij elke kopie daarna door \omterm{def:g11:mutations:mutation}{mutatie} naar
een eigen functie uiteengroeit (\omterm{def:g11:the-eye:photoreceptors}{opsines}, globines).
\end{solution}

\begin{solution}{exo:g12:diversification-of-life:3}
Haar \omterm{def:g10:universal-dna:information}{chromosomen}, één set van elke oudersoort, hebben bij de \omterm{def:g12:meiosis-genetic-shuffling:meiosis}{meiose} geen
homologen om mee te paren, zodat de gameten uit balans zijn. Verdubbeling
geeft elk \omterm{def:g11:cell-cycle-mitosis:chromosome}{chromosoom} een gelijke partner: de paring en de \omterm{def:g12:meiosis-genetic-shuffling:meiosis}{meiose} worden
regelmatig en de plant is vruchtbaar.
\end{solution}

\begin{solution}{exo:g12:diversification-of-life:4}
Hun eigen ringvormige \omterm{def:g10:universal-dna:information}{DNA}; \omterm{def:g10:cells-common-unit:organelle}{ribosomen} van het bacteriële type, gevoelig
voor antibacteriële \omterm{def:g11:antibiotic-resistance:antibiotic}{antibiotica}; vermenigvuldiging door deling; een dubbel
membraan; volgordes die het dichtst bij een groep levende bacteriën
liggen.
\end{solution}

\begin{solution}{exo:g12:diversification-of-life:5}
Een gedrag dat van andere leden van de groep wordt geleerd en zonder \omterm{def:g10:universal-dna:gene}{genen}
wordt doorgegeven: noten kraken met stenen in sommige
chimpanseegemeenschappen en niet in andere.
\end{solution}

\begin{solution}{exo:g12:diversification-of-life:6}
$14 + 7 = 21$ \omterm{def:g10:universal-dna:information}{chromosomen}: 7 A uit eenkoorn, 7 A en 7 B uit emmer. Er
kunnen zeven paren (A met A) worden gevormd; de 7 \omterm{def:g10:universal-dna:information}{B-chromosomen} hebben
geen partner en worden willekeurig verdeeld, zodat de gameten uit balans
zijn: onvruchtbaar.
\end{solution}

\begin{solution}{exo:g12:diversification-of-life:7}
De splitsing $\alpha$--$\beta$ (50\% gelijk) is ouder dan de splitsing
$\beta$--$\gamma$ (80\%): $\alpha$ takt eerst af, en daarna gaan $\beta$
en $\gamma$ uiteen.
\end{solution}

\begin{solution}{exo:g12:diversification-of-life:8}
De pootgenen zijn aanwezig, maar het signaal dat hen in de flank van het
embryo aanzet wordt niet gegeven, of wordt gegeven en daarna gestopt; het
verschil zit in de regeling van de expressie, niet in de \omterm{def:g10:universal-dna:gene}{genen}.
\end{solution}

\begin{solution}{exo:g12:diversification-of-life:9}
In de hoeveelheid en de timing van een groeisignaal in de embryonale
snavel: een verschil in regeling. Een verandering in de regeling vergt
slechts een kleine \omterm{def:g11:mutations:mutation}{mutatie} in een stuurvolgorde, terwijl een nieuw \omterm{def:g11:gene-expression:protein}{eiwit}
er vele zou vergen; zij kan binnen duizenden generaties ontstaan en worden
geselecteerd.
\end{solution}

\begin{solution}{exo:g12:diversification-of-life:10}
Zijn vorm, zijn vermogen om kaal gesteente te koloniseren, zijn weerstand
tegen droogte en zijn trage groei horen bij geen van beide partners
alleen; zij komen uit het verbond voort. Het korstmos is een samengesteld
organisme met eigen eigenschappen.
\end{solution}

\begin{solution}{exo:g12:diversification-of-life:11}
Doorgegeven gedrag: het materiaal is aan beide kanten beschikbaar, de
populaties zijn nauw verwant, en het verschil volgt de rivier --- een
belemmering voor contact en dus voor leren, niet voor \omterm{def:g10:universal-dna:gene}{genen} of
hulpbronnen.
\end{solution}

\begin{solution}{exo:g12:diversification-of-life:12}
Een virus bracht zijn \omterm{def:g10:universal-dna:gene}{gen} enkele tientallen miljoenen jaren geleden in het
genoom van een geslachtscel van een voorouder van de primaten; het \omterm{def:g10:universal-dna:gene}{gen}
werd geërfd en behouden. Toets: hetzelfde \omterm{def:g10:universal-dna:gene}{gen} zou bij alle primaten op
dezelfde plaats op het \omterm{def:g11:cell-cycle-mitosis:chromosome}{chromosoom} moeten zitten en op die plaats bij
andere zoogdieren moeten ontbreken; zijn volgorde zou het dichtst bij die
van het virus moeten liggen.
\end{solution}

\begin{solution}{exo:g12:diversification-of-life:13}
De \omterm{def:g10:cells-common-unit:organelle}{ribosomen} van de \omterm{def:g10:cells-common-unit:organelle}{mitochondriën} zijn van het bacteriële type en binden
dezelfde \omterm{def:g11:antibiotic-resistance:antibiotic}{antibiotica}: bij hoge doses vertraagt het middel de
eiwitsynthese van de \omterm{def:g10:cells-common-unit:organelle}{mitochondriën}. Die gevoeligheid is van de bacteriële
voorouder van de \omterm{def:g10:cells-common-unit:organelle}{mitochondriën} geërfd.
\end{solution}

\begin{solution}{exo:g12:diversification-of-life:14}
Het aantal \omterm{def:g10:universal-dna:information}{chromosomen} van de \omterm{prop:g12:diversification-of-life:polyploidy}{polyploïde} komt niet meer overeen met dat
van beide ouders, dus geven kruisingen met hen onvruchtbare nakomelingen:
de afzondering is onmiddellijk. Een nieuw \omterm{def:g10:universal-dna:gene}{allel} moet zich over vele
generaties door een populatie verbreiden voordat het een groep
onderscheidt. Een plantensoort kan dus in één enkele generatie ontstaan.
\end{solution}

\begin{solution}{exo:g12:diversification-of-life:15}
In de grond is elke nieuwe volgorde --- een verdubbelde kopie die
uiteengroeit, een verandering in de regeling, de \omterm{def:g10:universal-dna:gene}{genen} van een
overgedragen plasmide --- ergens als \omterm{def:g11:mutations:mutation}{mutatie} begonnen. Maar verdubbeling,
overdracht en polyploïdie verplaatsen en vermenigvuldigen hele \omterm{def:g10:universal-dna:gene}{genen} en
genomen in één keer; een verandering in de regeling brengt nieuwe vormen
voort zonder nieuwe \omterm{def:g10:chemistry-of-life:families}{eiwitten}; \omterm{def:g12:diversification-of-life:symbiosis}{symbiose} combineert genomen van
verschillende \omterm{def:g10:biodiversity-scales:species}{soorten}; en cultuur draagt variatie over zonder enig \omterm{def:g10:universal-dna:gene}{gen}.
De \omterm{def:g11:mutations:mutation}{mutatie} levert de letters; deze processen herschikken hele bladzijden
en boeken.
\end{solution}

\begin{solution}{pb:g12:diversification-of-life:1}
\textbf{1.} Emmer 28, broodtarwe 42.

\textbf{2.} Geen: de A- en \omterm{def:g10:universal-dna:information}{B-chromosomen} zijn niet \omterm{def:g10:common-ancestry:homology}{homoloog}. De gameten
krijgen willekeurige verzamelingen van de 14, vrijwel nooit een volledige
set.

\textbf{3.} Met 28 \omterm{def:g10:universal-dna:information}{chromosomen} heeft elke A een gelijke A-partner en elke
B een B: 14 paren, een regelmatige \omterm{def:g12:meiosis-genetic-shuffling:meiosis}{meiose}, gameten van 14 in balans.

\textbf{4.} Geen enkele kan paren: 7 A-, 7 B- en 7 \omterm{def:g10:universal-dna:information}{D-chromosomen}, elk
zonder \omterm{def:g10:common-ancestry:homology}{homoloog}. Onvruchtbaar.

\textbf{5.} 21 paren; een stuifmeelkorrel draagt er 21.

\textbf{6.} De andere twee kopieën leveren het \omterm{def:g11:gene-expression:protein}{eiwit} nog steeds.
Overtollige kopieën zijn vrij om \omterm{def:g11:mutations:mutation}{mutaties} op te hopen en uiteen te groeien
--- verdubbeling op de schaal van een heel genoom.

\textbf{7.} Op akkers in het gebied waar emmer werd verbouwd, tussen
10\,000 en 8\,000 jaar geleden, met wild geitengras 2 als onkruid tussen
het gewas.

\textbf{8.} Dat de twee stappen --- kruising en verdubbeling ---
volstaan om broodtarwe uit haar ouders voort te brengen: de geschiedenis
is opnieuw af te draaien.

\textbf{9.} Ja: zij kan met eenkoorn geen vruchtbare nakomelingen
voortbrengen, dus is zij volgens de definitie een aparte \omterm{def:g10:biodiversity-scales:species}{soort}, ook al
stamt zij ervan af.

\textbf{10.} De \omterm{def:g10:universal-dna:information}{D-chromosomen} van geitengras 2 zijn \omterm{def:g10:common-ancestry:homology}{homoloog} met de D-set
van de tarwe en paren ermee, zodat een kruising gevolgd door terugkruisen
het \omterm{def:g10:universal-dna:gene}{gen} kan overbrengen; de \omterm{def:g10:universal-dna:information}{chromosomen} van een niet-verwante plant zouden
helemaal niet paren.

\textbf{11.} De \omterm{def:g10:universal-dna:gene}{genen} zelf verschillen nauwelijks; het verschil in de
lichamen moet dus in het gebruik van de \omterm{def:g10:universal-dna:gene}{genen} zitten.

\textbf{12.} Een verandering in de regeling van de expressie --- de
uitgestrektheid van het gebied langs de as waar het \omterm{def:g10:universal-dna:gene}{gen} actief is.

\textbf{13.} De pootgenen zijn aanwezig en beginnen te werken; het signaal
dat de groei van de poot volhoudt ontbreekt, dus houden de knopjes op.
Opnieuw de regeling, en niet het verlies van de \omterm{def:g10:universal-dna:gene}{genen}.

\textbf{14.} Zij laat zien dat het veranderen van waar het \omterm{def:g10:universal-dna:gene}{gen} tot uiting
komt volstaat om het lichaamsgebied te veranderen: het verband tussen
expressie en vorm is oorzakelijk.

\textbf{15.} De verandering van broodtarwe voegde hele sets aan het genoom
toe (de omvang ervan); die van de python veranderde waar gedeelde \omterm{def:g10:universal-dna:gene}{genen}
worden gebruikt (hun regeling).

\textbf{16.} Eigen ringvormig \omterm{def:g10:universal-dna:information}{DNA} (verwacht: geen, alle \omterm{def:g10:universal-dna:gene}{genen} in de kern);
bacteriële \omterm{def:g10:cells-common-unit:organelle}{ribosomen} (verwacht: eukaryote \omterm{def:g10:cells-common-unit:organelle}{ribosomen}); vermenigvuldiging
door deling (verwacht: opbouw uit onderdelen); een dubbel membraan
(verwacht: één membraan zoals bij andere \omterm{def:g10:cells-common-unit:organelle}{organellen}). En hun \omterm{def:g10:universal-dna:gene}{genen} lijken
op bacteriële \omterm{def:g10:universal-dna:gene}{genen} (verwacht: \omterm{def:g10:universal-dna:gene}{genen} van het kerntype).

\textbf{17.} In de \omterm{def:g10:cells-common-unit:organelle}{celkern}: in de loop van de tijd zijn de meeste \omterm{def:g10:universal-dna:gene}{genen}
van de bacterie naar het kerngenoom overgebracht, en is het eigen genoom
van het \omterm{def:g10:cells-common-unit:organelle}{organel} tot een restant geslonken.

\textbf{18.} Mitochondriale \omterm{def:g10:universal-dna:gene}{genen} gaan alleen van moeder op kind over; het
mitochondriale \omterm{def:g10:universal-dna:information}{DNA} van iemand volgt de ononderbroken moederlijke lijn.

\textbf{19.} Het \omterm{def:g10:cells-common-unit:organelle}{mitochondrion} heeft de meeste van zijn \omterm{def:g10:universal-dna:gene}{genen} aan zijn
gastheer afgestaan en kan buiten de \omterm{def:g10:cells-common-unit:cell}{cel} niet bestaan; de partners van het
korstmos behouden hun genoom en zijn, met moeite, te scheiden.

\textbf{20.} Eenkoorn (14, A), geitengras 1 (14, B), geitengras 2 (14, D):
42 \omterm{def:g10:universal-dna:information}{chromosomen}; "regeling"; het \omterm{def:g10:cells-common-unit:organelle}{mitochondrion}.
\end{solution}
